\chapter{\MakeUppercase{Conclusion}}
\setcounter{equation}{0}

Stroke rehabilitation aims to improve motor function and independence, but patients often respond differently to the same treatment. Brain--Computer Interface (BCI)-based rehabilitation offers a promising approach to supporting motor recovery by connecting brain activity with assisted movement. However, predicting which patients will benefit from BCI treatment remains an important challenge.

This seminar reviewed the research paper ``Multimodal Behavioral Data Predict Stroke Patient's Response to BCI Treatment Through Explainable AI.'' The study combines demographic, clinical, and biomechanical characteristics to predict post-treatment lower-limb motor function and classify patients as responders or non-responders. It employs univariate and multivariate regression, Elastic Net, and Artificial Neural Network models, along with stepwise regression and SHapley Additive exPlanations (SHAP) to interpret model predictions.\cite{bib1}

The reported results demonstrate that the multivariate regression model achieved an adjusted $R^2$ of 0.852, compared with 0.758 for the univariate model. Elastic Net achieved a classification accuracy of 81.0\%, outperforming ANN at 76.2\%. Baseline motor function, muscle spasticity, and balance function were identified as important predictors of rehabilitation outcomes.

The study highlights the value of combining multimodal patient information with explainable machine learning. By identifying characteristics associated with predicted recovery, such frameworks may support more individualized and transparent rehabilitation planning. Nevertheless, the small sample size, limited treatment duration, and need for independent validation must be considered before applying the findings in routine clinical practice.

In conclusion, Explainable Artificial Intelligence provides a promising direction for understanding and predicting individual responses to BCI-based stroke rehabilitation. Further research using larger datasets, longer follow-up periods, and clinical validation is essential to establish practical effectiveness and support reliable, personalized rehabilitation strategies.